\chapter{Conclusion} \label{cap:conclusion}

The final result shows that all three criteria killed the same set of \texttt{longestPrefixOf} mutants.
Even with identical final scores, the criteria provide complementary evidence: edge-pair tests emphasize local transition correctness, prime-path tests stress deeper path behavior, and CDC targets decision sensitivity.
Using these perspectives together increases confidence in both control-flow and logic-level correctness.

For completeness, the remaining non-\texttt{longestPrefixOf} surviving mutants are boundary changes in
\texttt{collect} (lines 219 and 229), used by \texttt{keysThatMatch}.
They are equivalent, since changing \texttt{<} to \texttt{<=} and \texttt{>} to \texttt{>=} at those points the nodes exist, and the extra explored branch is always null and produces no output.

PIT's mutation analysis provided detailed insights into which test obligations were most effective at killing mutants, reinforcing the value of using multiple criteria to achieve comprehensive test coverage and fault detection.

Accompanied with logic-based criterion tests and edge-pair, prime-path and CDC tests, the test suit is now robust against a wide range of faults.